\usepackage[
  left=2cm,
  right=2cm,
  top=2.5cm,
  bottom=2.5cm,
  headheight=15pt,
  footskip=1.2cm
]{geometry}
\usepackage{fontspec}
\usepackage{setspace}
\usepackage{graphicx}
\usepackage{caption}
\usepackage{booktabs}
\usepackage{longtable}
\usepackage{array}
\usepackage{tabularx}
\usepackage{amsmath}
\usepackage{amssymb}
\usepackage{fancyhdr}
\usepackage{lastpage}
\usepackage{enumitem}
\usepackage{float}
\usepackage[hidelinks]{hyperref}

% 中文字体回退链：
%   1. WSL 下的 Windows 宋体/黑体（与已提交 PDF 一致）
%   2. 系统已安装的 SimSun/SimHei（原生 Windows/macOS）
%   3. 思源宋体/黑体 Noto Serif/Sans CJK SC（云端 Ubuntu，apt 包 fonts-noto-cjk）
%   4. TeX Live 自带的 Fandol
\IfFileExists{/mnt/c/Windows/Fonts/simsun.ttc}{%
  \setCJKmainfont[
    Path=/mnt/c/Windows/Fonts/,
    BoldFont=simhei.ttf
  ]{simsun.ttc}
  \setCJKsansfont[
    Path=/mnt/c/Windows/Fonts/,
    AutoFakeBold=2
  ]{simhei.ttf}
  \setCJKfamilyfont{zhhei}[
    Path=/mnt/c/Windows/Fonts/,
    AutoFakeBold=2
  ]{simhei.ttf}
}{\IfFontExistsTF{SimSun}{%
  \setCJKmainfont[BoldFont=SimHei]{SimSun}
  \setCJKsansfont[AutoFakeBold=2]{SimHei}
  \setCJKfamilyfont{zhhei}[AutoFakeBold=2]{SimHei}
}{\IfFileExists{/usr/share/fonts/opentype/noto/NotoSerifCJK-Regular.ttc}{%
  % FontIndex=2 为 ttc 中的简体中文（SC）字面，避免误用日文字形
  \defaultCJKfontfeatures{Path=/usr/share/fonts/opentype/noto/,Extension=.ttc,FontIndex=2}
  \setCJKmainfont[BoldFont=NotoSansCJK-Bold,BoldFeatures={FontIndex=2}]{NotoSerifCJK-Regular}
  \setCJKsansfont[BoldFont=NotoSansCJK-Bold,BoldFeatures={FontIndex=2}]{NotoSansCJK-Regular}
  \setCJKfamilyfont{zhhei}[BoldFont=NotoSansCJK-Bold,BoldFeatures={FontIndex=2}]{NotoSansCJK-Regular}
  \defaultCJKfontfeatures{}
}{%
  \setCJKmainfont{FandolSong-Regular}
  \setCJKsansfont{FandolHei-Regular}
  \setCJKfamilyfont{zhhei}{FandolHei-Regular}
}}}
\providecommand{\heiti}{\CJKfamily{zhhei}}
\setmainfont{TeX Gyre Termes}
\setsansfont{TeX Gyre Heros}

\setstretch{1.2}
\setlength{\parindent}{2em}
\setlength{\parskip}{0pt}
\setlength{\textfloatsep}{8pt plus 2pt minus 3pt}
\setlength{\floatsep}{8pt plus 2pt minus 3pt}
\setlength{\intextsep}{8pt plus 2pt minus 3pt}

\ctexset{
  section = {
    format = \zihao{3}\heiti,
    indent = 0pt,
    beforeskip = 1.0\baselineskip,
    afterskip = 0.5\baselineskip
  },
  subsection = {
    format = \zihao{4}\heiti,
    indent = 0pt,
    beforeskip = 0.8\baselineskip,
    afterskip = 0.35\baselineskip
  },
  subsubsection = {
    format = \zihao{-4}\heiti,
    indent = 0pt,
    beforeskip = 0.6\baselineskip,
    afterskip = 0.25\baselineskip
  }
}

\setlist[enumerate]{leftmargin=2.8em,itemsep=0.25em,topsep=0.25em}
\setlist[itemize]{leftmargin=2.5em,itemsep=0.25em,topsep=0.25em}
\captionsetup{font=small,labelsep=quad}
\renewcommand{\arraystretch}{1.25}
\newcolumntype{Y}{>{\centering\arraybackslash}X}

\fancypagestyle{report}{%
  \fancyhf{}
  \renewcommand{\headrulewidth}{0pt}
  \fancyfoot[C]{\zihao{-5}浙江大学物理实验教学中心制\quad 第~\thepage~页~/~共~\pageref{LastPage}~页}
}
\fancypagestyle{plain}{%
  \fancyhf{}
  \renewcommand{\headrulewidth}{0pt}
  \fancyfoot[C]{\zihao{-5}浙江大学物理实验教学中心制\quad 第~\thepage~页~/~共~\pageref{LastPage}~页}
}

\newcommand{\InfoField}[2][7cm]{\underline{\makebox[#1][c]{#2}}}
